\section{Deep Exploration: From MDPs to LLMs}

\begin{frame}{}
    \LARGE Exploration: \textbf{From Bandits to Deep RL}
\end{frame}

\begin{frame}{Why Bandit Exploration Doesn't Scale Naively}
\begin{itemize}
    \setlength{\itemsep}{0.55em}
    \item Bandits gave us \emph{optimism} (UCB) and \emph{posterior sampling} (Thompson) --- both in a single step, with no state to worry about
    \item In an MDP with sparse rewards, none of these scale naively:
    \begin{itemize}
        \item Visiting every $(s,a)$ enough times is infeasible
        \item $\epsilon$-greedy is provably weak (linear regret)
        \item Uniform random exploration drowns in pixels
    \end{itemize}
    \item Note: SAC's $-\log\pi$ entropy bonus (Day 6) is \emph{undirected} exploration --- what follows is \emph{directed}
\end{itemize}
\end{frame}

\begin{frame}{The Through-Line: Montezuma's Revenge}
\begin{center}
\textit{The canonical hard-exploration Atari game.}
\end{center}
\vspace{0.6em}
\begin{itemize}
    \setlength{\itemsep}{0.45em}
    \item \textbf{Vanilla DQN:} $\sim 0$ points (extrinsic reward is essentially never reached)
    \item \textbf{Pseudo-counts (2016):} first dramatic progress
    \item \textbf{Random Network Distillation (2018):} first above-human, no demonstrations
    \item \textbf{Go-Explore (2019, Nature 2021):} $>2{,}000{,}000$ points --- solves \emph{all} levels
    \item \textbf{Agent57 (2020):} above human on \emph{all 57} Atari games, Montezuma included
\end{itemize}
\vspace{0.5em}
\centering\small\textcolor{gray}{This is the single benchmark we'll use to anchor every method in this section.}
\end{frame}

\begin{frame}{Aside: Hindsight Experience Replay (HER)}
\textbf{An orthogonal trick for sparse-reward goal-conditioned RL.}
\begin{itemize}
    \setlength{\itemsep}{0.55em}
    \item Idea: after a failed episode (didn't reach the goal), relabel it as a \emph{success} for the goal it actually reached
    \item Now every episode has a non-zero learning signal
    \item Combines with any off-policy algorithm
    \item Cite: \textbf{Andrychowicz et al., NeurIPS 2017} --- \emph{Hindsight Experience Replay}
\end{itemize}
\end{frame}

\begin{frame}{Bootstrapped DQN: Thompson Sampling for MDPs}
\begin{itemize}
    \setlength{\itemsep}{0.55em}
    \item Idea: \emph{Posterior Sampling for RL (PSRL).} Maintain a posterior over the Q-function; sample one Q per episode and act greedily w.r.t.\ the sample
    \item For an \emph{entire episode}: temporally-extended (``deep'') exploration --- fundamentally different from $\epsilon$-greedy's per-step dithering
\end{itemize}
\end{frame}

\begin{frame}{Bootstrapped DQN: How It Approximates the Posterior}
\begin{itemize}
    \setlength{\itemsep}{0.32em}
    \item \textbf{$K$ parallel Q-heads}: $K$ output layers on one shared trunk, each with its own Q-function
    \item \textcolor{gray}{``Replay'' is DQN's \textbf{replay buffer} from Day 2 --- the stored $(s,a,r,s')$ transitions so far}
    \item Each transition is assigned to only \emph{some} of the $K$ heads, so every head trains on a different slice of experience
    \item \emph{Each episode}: pick one head at random and act greedily by it throughout
    \item \textbf{Where data is thin the heads disagree; where it is plentiful they converge} --- that spread \emph{is} the posterior estimate
    \item \textcolor{red}{So disagreement is the product, not a defect to train away.} Heads that were all good everywhere would make sampling identical to greedy --- plain DQN, no exploration
    \item More coverage $\Rightarrow$ heads converge $\Rightarrow$ exploration fades by itself. \textcolor{gray}{Thompson sampling lifted to MDPs: $K$ point-samples instead of a curve (Osband et al., 2016)}
\end{itemize}
\end{frame}

\begin{frame}{Count-Based Exploration: The Idea}
\textbf{The shape every method here shares.} We do not change the algorithm --- we change the reward it optimises:
\[
    r_t^+ \;=\; \underbrace{r_t}_{\text{real reward from the env}} \;+\; \beta \cdot \underbrace{\frac{1}{\sqrt{N(s,a)}}}_{\text{novelty bonus}}
\]
\begin{itemize}
    \setlength{\itemsep}{0.4em}
    \item \textcolor{red}{The bonus is simply \emph{added} to the real reward}, and PPO/DQN/SAC run unchanged --- from here on, only the \textbf{bonus term} differs between methods
    \item A direct lift of UCB1, with $\beta$ trading exploration against the real task
    \item \emph{Problem once the state is an image} (Atari, a camera feed): the state is the \textbf{whole frame}, so two states match only if \emph{every pixel} does
    \item So $N(s,a) = 1$ forever and the bonus is constant everywhere. We need a \emph{generalizable} count --- one that calls two \emph{similar} frames both familiar
\end{itemize}
\end{frame}

\begin{frame}{Pseudo-Counts: UCB in Pixel Space}
\textbf{A density model} $\rho$ answers one question: \emph{``how likely is this frame?''} It is trained on the frames seen so far, so familiar-looking screens get high $\rho$ and never-seen ones get low $\rho$.
\vspace{0.3em}
\begin{itemize}
    \setlength{\itemsep}{0.45em}
    \item \textbf{The trick:} train the model on \emph{one more copy of $s$} and see how much its answer for $s$ moves.
    \[
        \hat{N}(s) \;=\; \frac{\rho(s)\,\big(1 - \rho'(s)\big)}{\rho'(s) - \rho(s)}
    \]
    \item \textcolor{red}{$\rho'(s)$ is not the next state.} It is the \emph{same} state $s$, scored by the model \emph{after} that extra training step --- ``before'' and ``after'', not ``now'' and ``next''
    \item \textbf{Why this is a count:} a model that has seen $s$ a million times barely budges, so the denominator is tiny and $\hat{N}$ is huge. A model seeing $s$ for the first time jumps, so $\hat{N}$ is small --- and the bonus is large
    \item Cite: \textbf{Bellemare et al., NeurIPS 2016}. \textbf{Result on Montezuma:} \emph{first dramatic progress}
\end{itemize}
\end{frame}

\begin{frame}{Curiosity / ICM: The Idea}
\textbf{Intrinsic reward $=$ ``surprise''} --- the agent is rewarded for finding observations it can't yet predict.
\begin{itemize}
    \setlength{\itemsep}{0.55em}
    \item Embed states into a feature space: $\phi(s) \in \mathbb{R}^d$
    \item Train a \textbf{forward-dynamics model} $f$ to predict the \emph{next-state features} from the current features and action:
    \[
        \hat{\phi}(s_{t+1}) \;=\; f\!\big(\phi(s_t),\, a_t\big)
    \]
    (so $\hat{\phi}(s_{t+1})$ is the model's \emph{predicted} embedding of $s_{t+1}$; $\phi(s_{t+1})$ is the actual one observed)
    \item The bonus is how wrong that prediction was --- \textcolor{red}{added to the real reward as before}, $r_t^+ = r_t + \beta\, r_t^i$:
    \[
        r_t^i \;=\; \tfrac{1}{2}\big\|\,\hat{\phi}(s_{t+1}) \,-\, \phi(s_{t+1})\,\big\|^2
    \]
    \item High prediction error $\Rightarrow$ novel transition $\Rightarrow$ bonus to visit it again
\end{itemize}
\end{frame}

\begin{frame}{Curiosity / ICM: Why the Feature Space Matters}
\textbf{The noisy-TV problem.} Picture a TV in the environment showing pure random static. A naive curiosity agent will sit and stare at it forever --- static is \emph{always} unpredictable, so the prediction-error bonus never decays. The agent is endlessly ``surprised'' by useless noise.
\vspace{0.4em}
\begin{itemize}
    \setlength{\itemsep}{0.5em}
    \item Fix: don't let $\phi$ encode action-irrelevant detail. Train $\phi$ with an \textbf{inverse-dynamics} objective --- predict the action $a_t$ that was taken from the embeddings $\big(\phi(s_t),\, \phi(s_{t+1})\big)$
    \item This forces $\phi$ to encode only features the agent can \emph{affect} --- the noisy TV is filtered out, since the agent's actions don't change what's on it
    \item Cite: \textbf{Pathak, Agrawal, Efros \& Darrell, ICML 2017} --- \emph{Curiosity-driven Exploration by Self-supervised Prediction}
    \item Strong results on Mario, VizDoom, and sparse-reward navigation
\end{itemize}
\end{frame}

\begin{frame}{Random Network Distillation (RND)}
\textbf{Curiosity made absurdly simple:}
\begin{itemize}
    \setlength{\itemsep}{0.55em}
    \item Pick a \emph{fixed, randomly-initialized} target network $f^*$
    \item Train a predictor $f_\theta$ to match $f^*(s)$ on visited states
    \item Bonus: $r_t^i = \|f_\theta(s_t) - f^*(s_t)\|^2$, again \textcolor{red}{added to the real reward}, $r_t^+ = r_t + \beta\, r_t^i$
\end{itemize}
\vspace{0.6em}
\begin{itemize}
    \item Novel states have high prediction error (the predictor hasn't been trained there) $\Rightarrow$ high bonus
    \item No dynamics model, no inverse model, no density model
\end{itemize}
\end{frame}

\begin{frame}{RND: Results and Caveat}
\begin{itemize}
    \setlength{\itemsep}{0.55em}
    \item Cite: \textbf{Burda, Edwards, Storkey, Klimov, ICLR 2019} --- \emph{Exploration by Random Network Distillation}
    \item \textbf{Result on Montezuma:} \emph{first above-human performance without demonstrations or game-state access}
    \item \emph{Caveat:} still subject to the noisy-TV issue in principle; in practice works because randomly-initialized networks aren't very sensitive to stochastic pixel noise
    \item \textcolor{red}{Still the default intrinsic-reward baseline in 2026} --- simple, cheap, and reliable. Recent work refines rather than replaces it (Distributional RND; Random Distribution Distillation, 2025)
\end{itemize}
\end{frame}

\begin{frame}{Go-Explore: Remember Where You Have Been}
\textbf{A different diagnosis.} The problem isn't only \emph{finding} novel states --- it's that agents \emph{forget} promising ones and can never get back to them.
\vspace{0.4em}
\begin{itemize}
    \setlength{\itemsep}{0.5em}
    \item Keep an \textbf{archive} of visited states. Each iteration: pick a promising one, \emph{return} to it directly, then explore \emph{from} there --- ``first return, then explore''
    \item Exploration effort is never wasted re-deriving a path you already found
    \item \textbf{Result:} $>2{,}000{,}000$ on Montezuma and all levels solved --- two orders of magnitude past RND
    \item \emph{The catch:} returning to an archived state originally needed simulator save/restore, a much stronger assumption than RND's. That is why both are ``state of the art'' in different senses
    \item Cite: \textbf{Ecoffet, Huizinga, Lehman, Stanley \& Clune} --- \emph{Go-Explore} (2019); \emph{First Return, Then Explore}, Nature 2021
\end{itemize}
\end{frame}

\begin{frame}{Episodic vs.\ Lifelong Novelty}
\textbf{A distinction the methods so far miss.} RND asks ``have I \emph{ever} seen this?'' But within a single episode you also want to avoid re-treading the same room --- even if you have seen it many times before.
\vspace{0.4em}
\begin{itemize}
    \setlength{\itemsep}{0.5em}
    \item \textbf{Never Give Up (NGU, 2020)} combines both signals: an \emph{episodic} novelty bonus (reset each episode) multiplied by a \emph{lifelong} one (RND-style, across all training)
    \item \textbf{Agent57 (2020)} adds a meta-controller that \emph{learns} how much to weight exploration versus exploitation, and over what horizon --- first agent above human on all 57 Atari games
    \item \textbf{E3B (2022)} is the modern count-based descendant: an episodic bonus for states that are \emph{diverse} under an inverse-dynamics embedding --- ICM's noisy-TV fix combined with UCB's count bonus, extended to continuous spaces
    \item \textcolor{red}{The pattern:} no single novelty signal is enough --- current systems \emph{combine} them at different timescales
\end{itemize}
\end{frame}

\begin{frame}{Exploration Did Not Go Away --- It Moved to LLMs}
\textbf{The same dilemma, today's frontier.} When you train a reasoning model with RL (Day 9), it faces exactly the trade-off we opened this morning with.
\vspace{0.4em}
\begin{itemize}
    \setlength{\itemsep}{0.5em}
    \item \textbf{Entropy collapse:} the policy sharpens onto a few reasoning paths early in training. In-distribution accuracy keeps rising while out-of-distribution performance \emph{degrades} --- pure exploitation, and the model stops discovering
    \item \textbf{How it is measured:} \emph{pass@k} --- can the model solve the problem in $k$ tries? Collapse shows up as pass@1 improving while pass@k falls. Training directly on pass@k restores exploration
    \item \textbf{Today's fixes are today's bandit ideas:} entropy bonuses aimed at the most uncertain tokens (Day 6's $-\log\pi$), and \textcolor{red}{count-based bonuses of the form $Q + \alpha\sqrt{U}$} --- literally UCB's $1/\sqrt{N}$, applied to reasoning states
\end{itemize}
\vspace{0.3em}
\centering\small\textcolor{gray}{The 1985 bandit question --- ``is this worth trying again?'' --- is an open problem in 2026 LLM training.}
\end{frame}

\begin{frame}{Mapping Bandits $\longrightarrow$ Deep RL $\longrightarrow$ LLMs}
\begin{center}
\renewcommand{\arraystretch}{1.4}
\footnotesize
\begin{tabular}{p{0.28\textwidth}p{0.60\textwidth}}
\hline
\textbf{Bandit idea} & \textbf{Descendants} \\
\hline
UCB / $1/\sqrt{N}$ bonus & Pseudo-counts (2016), E3B (2022); \textcolor{red}{count bonuses in LLM RL (2025)} \\
Thompson sampling     & Bootstrapped DQN / PSRL (Osband 2016) \\
Stochastic policies   & \textcolor{blue}{Entropy bonus --- SAC (Day 6)}; entropy control in RLVR \\
\emph{(new in deep RL)} & Curiosity / ICM (2017); RND (2018); NGU / Agent57 (2020); Go-Explore (2021) \\
\hline
\end{tabular}
\end{center}
\vspace{0.5em}
\begin{itemize}
    \setlength{\itemsep}{0.45em}
    \item \textcolor{red}{\textbf{Entropy is only one answer to exploration.}} Optimism, posterior sampling, novelty, and archiving are the others
    \item \textbf{What actually survived:} RND is still the default baseline; the $1/\sqrt{N}$ bonus keeps reappearing; $\epsilon$-greedy is still what most people try first
\end{itemize}
\end{frame}
